\section{Related Work}
\label{sec:related-work}

\subsection{Prediction Time and Leakage in ED Risk Models}

Leakage, meaning information in the model input that would not be available when the model is used, has been documented across machine-learning-based
science~\cite{kapoor2023leakage}. In MIMIC-IV, diagnostic codes that are finalized only after discharge inflate the performance of same-admission
models~\cite{ramadan2025icdleakage}. Avoiding label leakage takes more nuance than a rule against using the future: Davis et al.\ propose to consider the cadence, the perspective, and the applicability of a prediction~\cite{davis2023labelleakage}. The presence and timing of measurements can themselves carry information about a patient, which methods for clinical prediction models treat as informative presence~\cite{sisk2021informativepresence}, and the value of such a predictor depends on the transportability of the mechanism that generates the missing data~\cite{groenwold2020informativemissingness}. Shortcut learning describes models that rely on features that perform well on a benchmark and do not transfer to the intended
use~\cite{geirhos2020shortcut}. In survival analysis, landmarking restricts both the inputs and the at-risk population to what holds at a chosen
time~\cite{vanhouwelingen2007landmarking}. For ED outcomes, models built on triage data fix the prediction time by construction~\cite{raita2019triage,hong2018admission}. The
MIMIC-IV-ED benchmark~\cite{xie2022mimic4ed} predicts hospitalization and a critical outcome at triage and return to the ED at disposition, and in its comparison deep
sequence models did not outperform simpler ones; its pipeline also provides the vital signs charted during the stay, which makes predictions from the whole ED record
easy to build. A multimodal benchmark on the same database adds electrocardiogram waveforms and laboratory values to predict discharge diagnoses and deterioration
targets~\cite{alcaraz2025mdsed}, and a leakage-aware benchmark restricts the inputs to the first hour of presentation and compares vital-sign-only with richer
inputs~\cite{solaiman2026earlydeterioration}. CAREBench matches the length of the observed history between patients with and without the event, so that a
model cannot use it to separate them~\cite{keoliya2025carebench}. These designs prevent a timing shortcut; they do not measure how large it is. In a recorded search (Supplementary Table~\ref{S-tab:search}), none of \RVLitHitsOne{} records on ED prediction models with workflow
shortcuts or leakage, of the first 60 of \RVLitHitsSix{} records on informative presence, monitoring frequency, or label leakage in clinical prediction, or of \RVLitHitsSeven{} records on ED models with a stated prediction time quantified how much of the discrimination of an end-of-stay model comes from quantities that describe how a stay was monitored; the closest work concerns diagnostic codes~\cite{ramadan2025icdleakage} and informative presence in general~\cite{sisk2021informativepresence}.

\subsection{Temporal Models for Vital Signs}

Deep early warning systems learn from vital-sign streams to detect clinical deterioration~\cite{shamout2020dews}. GRU-D handles irregular sampling by decaying each
variable toward its last observed value as a function of elapsed time~\cite{che2018grud}, and set functions~\cite{horn2020seft} and multi-time attention
networks~\cite{shukla2021mtand} address the same setting; gradient-boosted trees nevertheless remain strong on tabular data~\cite{grinsztajn2022treebased}. Clinical
models lose discrimination under temporal dataset shift~\cite{nestor2019featurerobustness,guo2022domaingeneralization}, so deployed models need continuing
evaluation~\cite{kim2025monitoring}. CAREBench also evaluates the temporal stability of the risk score across electronic health records, waveforms, and clinical
text, and reports poor recall at high-precision operating points~\cite{keoliya2025carebench}. These studies measure how discrimination or the stability of the score changes, not whether the uncertainty attached to a risk score stays valid on a later period.

\subsection{Conformal Prediction for Clinical Risk}

Split conformal prediction gives finite-sample marginal coverage under exchangeability~\cite{lei2018distributionfree,vovk2022algorithmic}; for classification we use the
least-ambiguous set-valued score~\cite{sadinle2019lac} as implemented in MAPIE~\cite{cordier2023mapie}. Clinical applications have calibrated sepsis mortality
risk~\cite{yang2024sepsisconformal}, sepsis diagnosis with external validation~\cite{dalal2025sepsisconformal}, and early sepsis prediction with cost-aware deferral to
clinicians~\cite{kwon2026conformaltriage}, and uncertainty-aware models have been applied to disposition at ED triage~\cite{abdulai2025idontknow}. These studies use static
calibration, evaluated on a random split, an external cohort, or later-admitted patients. A scoping review found clinical conformal studies conducted in isolation~\cite{vazquez2022clinicalcp}, an overview surveyed conformal prediction in clinical artificial intelligence~\cite{wang2026conformalclinicalai}, and a critical perspective questioned the value of the finite-sample guarantees in medicine~\cite{kladny2025criticalperspective}. Marginal coverage does not imply coverage within a minority class; Mondrian conformal prediction calibrates within
categories such as the observed class~\cite{vovk2003mondrian,ding2023classconditional}, and a recent preprint on multiple sclerosis imaging found that marginal sets
under-covered the minority class under distribution shift while class-conditional calibration reduced the gap~\cite{millar2026classconditionalms}. For shift over time,
adaptive conformal inference updates the miscoverage level after each observation~\cite{gibbs2021aci}, with later work on step sizes and
stability~\cite{zaffran2022adaptivecp,angelopoulos2023pid}; weighted conformal prediction handles covariate shift~\cite{tibshirani2019covariateshift}; Barber et al.\ bound
the coverage gap without exchangeability~\cite{barber2023beyondexchangeability}; and ECI adds an error-quantification term to the adaptive update~\cite{wu2025eci}.
Delayed and intermittent feedback have been studied~\cite{elhalabi2026delayed,wang2025mirrorocp}. Among the screened records (\RVLitHitsTwo{} on conformal prediction
with ED, triage, deterioration or intensive care; \RVLitHitsThree{} on class-conditional conformal prediction in clinical data; \RVLitHitsFour{} on adaptive or online
conformal prediction under shift in clinical data; \RVLitHitsFive{} on delayed feedback), none reported the coverage of an ED deterioration model separately for the
patients who deteriorate together with the false-alarm burden of the sets, and none compared online methods with static and periodically recalibrated calibration on a
clinical stream; the delayed-feedback studies we identified were not clinical.

\subsection{Early Warning Scores}

NEWS2 is the UK standard early warning score~\cite{rcp2017news2} and has been examined in hospital and emergency settings~\cite{gerry2020ews}; in one ED cohort it reached
AUROCs of \RVRefNewsLo{} to \RVRefNewsHi{} across outcomes~\cite{mhdyunin2025news2}. Learned alternatives use the same vital-sign
streams~\cite{shamout2020dews,ibrahim2022kdensemble}. These studies evaluate discrimination and interpretability, not the coverage of prediction sets under temporal shift.

\subsection{Position of This Study}

Supplementary Table~\ref{S-tab:related} compares the closest studies with this one on the data, the prediction time, the handling of input timing, the
uncertainty method, and the evaluation under shift. Relative to the ED benchmarks~\cite{xie2022mimic4ed,alcaraz2025mdsed,solaiman2026earlydeterioration}, this study audits each input of an end-of-stay model against its
source timestamp and measures what quantities that describe how the stay was monitored add to arrival information. Relative to clinical conformal
studies~\cite{yang2024sepsisconformal,dalal2025sepsisconformal,kwon2026conformaltriage}, it evaluates the sets on a chronological split for the class that matters
clinically and reports their false-alarm burden across coverage levels. Relative to methods for conformal prediction under shift~\cite{gibbs2021aci,wu2025eci}, it compares
an adaptive method with static and periodic recalibration on a clinical stream with label delay. It does not propose a new prediction or conformal method.
